\chapter*{摘要}
\addcontentsline{toc}{chapter}{摘要}
\markboth{摘要}{摘要}

本书以“对称性 $\to$ 不变量 $\to$ 闭合轨道 $\to$ 周期”为纲领，构建了一套从第一原理出发统一周期现象的数学框架。

问题的起点是：不同学科中的周期现象（天体运行、电路振荡、生物节律、经济波动、算法迭代）长期使用互不相通的描述语言。本书主张，它们可以被同一个数学对象刻画——\textbf{周转圆族}：一个周期为 $T$ 的系统由振幅、角频率、相位三元组 $(A_k,\omega_k,\phi_k)$ 的集合描述，而周期之间的关系由\textbf{周期卷积}刻画。全书由此把“周期性”归约为群作用与群表示的问题：平移、缩放、旋转三种基本变换生成全部初等周期结构，而其完备性在复平面单参数保角变换族的框架内可被严格证明。

全书的核心方法论由四大支柱构成：群表示论、周期卷积、窗函数与折射率体系、复指数激活函数（CEAF）与核平滑导-积下降法。在此之上，本书给出把任意目标接入该框架的五步标准接口，并在天文学、物理学、生物学、心理学、经济学、密码学与计算机科学等二十余个领域中逐一套用与验证。书中严格区分三件事：已证明的定理、以研究纲领形式提出的构想、以及尚未解决的问题。

\vspace{1em}
\noindent\textbf{关键词：} 周期数学；周转圆；调和分析；守恒律；对称性；相位；振幅；频率；周期卷积；群表示论；傅里叶变换；梅林变换；小波变换；易经；莱布尼茨；历法；芯片；信号处理

\chapter*{Abstract}
\addcontentsline{toc}{chapter}{Abstract}
\markboth{Abstract}{Abstract}

This book develops a unified mathematical framework for periodic phenomena, based on the principle ``symmetry $\to$ invariant $\to$ closed orbit $\to$ period''.

The starting point is that periodic phenomena in different disciplines---celestial motion, electrical oscillation, biological rhythm, economic fluctuation and algorithmic iteration---have long been described in mutually untranslatable languages. The book argues that they can all be captured by a single mathematical object: a \textbf{family of epicycles}. A system of period $T$ is described by a set of amplitude, angular-frequency and phase triples $(A_k,\omega_k,\phi_k)$, whilst the relations between periods are captured by \textbf{periodic convolution}. Periodicity is thereby reduced to a question about group actions and group representations: the three basic transformations---translation, scaling and rotation---generate every elementary periodic structure, and their completeness can be proved rigorously within the framework of one-parameter families of conformal transformations of the complex plane.

The core methodology consists of four pillars: group representation theory, periodic convolution, window functions and refraction-index diagnostics, and complex exponential activation functions (CEAF) with kernel-smoothed derivative--integral descent. On this basis the book gives a standard five-step interface for attaching an arbitrary target to the framework, and applies it domain by domain in more than twenty fields, from astronomy, physics, biology and psychology to economics, cryptography and computer science. Throughout, three things are kept strictly apart: theorems that have been proved, conceptions advanced as research programmes, and problems that remain open.

\vspace{1em}
\noindent\textbf{Keywords:} periodic mathematics; epicycle; harmonic analysis; conservation law; symmetry; phase; amplitude; frequency; periodic convolution; group theory; Fourier transform; Mellin transform; wavelet transform; I Ching; Leibniz; calendar; chip; signal processing

\newpage

